\subsection{state}\label{state}
The structure \verb|state_t| is the central data structure of the system: it
holds the key state that is exchanged between all input devices (Section
\ref{keyboard}) and synchronised over the radio between the devices. Each of
the up to \verb|MAX_BUTTON_CNT| keys has a state \verb|key_state_e| --
\verb|OFF|, \verb|BLINKING| or \verb|ON| (Listing \ref{key_state_e}) -- as
well as a label.

The central mechanism is \verb|state_propagate()|: a key press advances the
state cyclically
($\mathtt{OFF} \to \mathtt{BLINKING} \to \mathtt{ON} \to \mathtt{OFF}$, i.e.\
$\mathtt{state} = (\mathtt{state}+1) \bmod \mathtt{STATE\_CNT}$) and marks the
\verb|state_t| as \emph{dirty}.

\begin{listing}
\begin{minted}{C}
typedef enum __attribute__((packed)) {
    OFF, BLINKING, ON, STATE_CNT
} key_state_e;
\end{minted}
\caption{Key state \texttt{key\_state\_e}}
\label{key_state_e}
\end{listing}

\subsubsection*{State}
Listing \ref{state_t} shows the structure (size $40$~bytes). \verb|first| and
\verb|cnt| define the active window of keys, \verb|dirty| marks a change since
the last synchronisation, \verb|id| is a counter carried along on every
transmission and \verb|clabel| (Section \ref{common}) carries a command or
label. The fields \verb|state[]| and \verb|label[]| hold state and label per
key.

\begin{listing}
\begin{minted}{C}
typedef struct state_s {
    uint8_t  first;                 // first valid key
    uint8_t  cnt;                   // number of valid keys
    uint8_t  dirty;                 // change flag
    uint8_t  id;                    // transmit counter
    clabel_u clabel;                // command/label (4 bytes)
    key_state_e state[MAX_STATE_CNT]; // state per key
    char     label[MAX_STATE_CNT];  // label per key
} state_t;                          // 2*MAX_BUTTON_CNT + 8 = 40 bytes
\end{minted}
\caption{Definition of the \texttt{state\_t} structure}
\label{state_t}
\end{listing}

\noindent Via the option \verb|REDUCED_PAYLOAD| (in \verb|emLibC_options.h|) a
bit-packed variant ($\approx 10$~bytes) exists alternatively, which is
disabled here.

\subsubsection*{Life cycle and individual keys}
{\sloppy\emergencystretch=3em
\begin{description}
\item[\texttt{state\_init(state)} / \texttt{state\_reset(state)}] Initialises
  the structure (default labels \verb|0..F|) or sets all keys to \verb|OFF|.
  Returns: \verb|em_msg| (\verb|EM_ERR| on an invalid structure, otherwise
  \verb|EM_OK|).
\item[\texttt{state\_check(state)}] Checks the range fields for validity.
  Returns: \verb|em_msg| (\verb|EM_OK| for valid range fields, otherwise
  \verb|EM_ERR|).
\item[\texttt{state\_set(state, nr, ns)} / \texttt{state\_get(state, nr)}]
  Sets or reads the state of a key by its index. Returns: \verb|em_msg|
  (setter; \verb|EM_ERR| on an invalid index/structure) or \verb|key_state_e|
  (the key state).
\item[\texttt{state\_set\_key\_by\_lbl(\dots)} /
  \texttt{state\_get\_key\_by\_lbl(\dots)}] The same via the label;
  \verb|state_ch2idx()| maps a character to the index (\verb|int8_t|, $-1$ if
  not found). Returns: \verb|em_msg| (setter) or \verb|key_state_e| (the key
  state, \verb|STATE_CNT| if not found).
\item[\texttt{state\_propagate(state, idx)} /
  \texttt{state\_propagate\_by\_idx} / \texttt{\dots\_by\_lbl}] Advances a key
  cyclically by one step and sets \emph{dirty}. Returns: \verb|em_msg|
  (\verb|EM_ERR| on an invalid structure, otherwise \verb|EM_OK|).
\end{description}
\par}

\subsubsection*{Whole state: transferring and combining}
{\sloppy\emergencystretch=3em
\begin{description}
\item[\texttt{state\_set\_state(from, to)} / \texttt{state\_copy(from, to)}]
  Copy the key content from \verb|from| to \verb|to|. Returns: \verb|em_msg|
  (\verb|EM_ERR| on an invalid structure, otherwise \verb|EM_OK|).
\item[\texttt{state\_set\_u32(state, u32)} / \texttt{state\_get\_u32(state)}]
  Pack or unpack the $16$ keys of $2$ bits each into a \verb|uint32_t|.
  Returns: \verb|em_msg| (setter) or \verb|uint32_t| (the packed states).
\item[\texttt{state\_diff(ref, state, diff)}] Determines the per-key difference
  (modular, cf.\ \verb|state_key_diff()|) relative to \verb|ref| in
  \verb|diff|. Returns: \verb|em_msg| (the \emph{dirty} flag).
\item[\texttt{state\_add(ref, add)}] Applies \verb|add| as a number of
  advances to \verb|ref| (\verb|BLINKING| $= +1$, \verb|ON| $= +2$). Returns:
  \verb|em_msg| (the \emph{dirty} flag).
\item[\texttt{state\_merge(in, out)}] Copies changed keys from \verb|in| to
  \verb|out| and reports whether anything changed. Returns: \verb|bool|.
\item[\texttt{state\_is\_same(last, cur)}] Checks two states for equality.
  Returns: \verb|em_msg| (\verb|EM_OK| if equal).
\end{description}
\par}

\subsubsection*{Range, flags and output}
{\sloppy\emergencystretch=3em
\begin{description}
\item[\texttt{state\_set\_first/cnt} / \texttt{state\_get\_first/cnt}] Set or
  read the active key window. Returns: \verb|em_msg| (setters) or
  \verb|uint8_t| (getters).
\item[\texttt{state\_get\_dirty} / \texttt{state\_set\_dirty} /
  \texttt{state\_set\_undirty}] Manage the change flag (\verb|state_get_dirty|
  returns its value). Returns: \verb|em_msg| (\verb|EM_ERR| on an invalid
  structure).
\item[\texttt{state\_print(state, title, doLong, cycle)}] Prints labels,
  states and the cycle position (Section \ref{cycle}) for debugging. Returns:
  \verb|em_msg| (\verb|EM_ERR| on an invalid structure, otherwise \verb|EM_OK|).
\end{description}
\par}
